\documentclass{book}

\title{Analisi Matematica T-A}
\author{Michele Mazza}
\date{November 2023}


% Pacchetti di Codifica e Lingua
\usepackage[utf8]{inputenc}
\usepackage[T1]{fontenc}
\usepackage[italian]{babel}

% Pacchetti Grafici e Diagrammi
\usepackage{graphicx}
\usepackage{tikz}
\usepackage{tikz-3dplot}
\usetikzlibrary{angles,quotes,arrows.meta, positioning}
\usepackage{pgfplots}
\pgfplotsset{compat=1.18}

% Pacchetti Matematici
\usepackage{amsmath, amssymb, amsthm}

% Pacchetti per Tabelle e Unità di Misura
\usepackage{geometry}
\usepackage{siunitx}
\usepackage{booktabs}
\usepackage{tabularx}
\usepackage{colortbl}

% Pacchetti per Layout e Formattazione
\usepackage{fancyhdr}
\usepackage{float}
\usepackage{microtype}

% Altri Pacchetti Utili
\usepackage{mdframed}
\usepackage{animate}

% Pacchetto enumitem
\usepackage{enumitem}

% Definizioni dei Colori
\definecolor{lightblue}{HTML}{9cd8fc}
\definecolor{myblue}{HTML}{139FE9}

% Definizione del comando \norm
\newcommand{\norm}[1]{\left\lVert#1\right\rVert}

%\font\nullfont=cmr10 

% Ambienti Personalizzati
\newmdtheoremenv[
    backgroundcolor=orange!10,
    linecolor=myblue,
    linewidth=0pt,
    leftmargin=0pt,
    rightmargin=0pt,
    innerleftmargin=10pt,
    innerrightmargin=10pt,
    innertopmargin=10pt,
    innerbottommargin=10pt,
    nobreak
]{curiosity}{Curiosità}

\newmdtheoremenv[
    font=\normalfont\small,
    linewidth=0pt,
    innerleftmargin=10pt,
    innerrightmargin=10pt,
]{example}{Esempio}

\newmdtheoremenv[
  backgroundcolor=lightblue,
  linecolor=myblue,
  linewidth=1.5pt,
  innertopmargin=10pt,
  innerbottommargin=10pt,
  innerrightmargin=10pt,
  innerleftmargin=10pt,
  nobreak
]{theorem}{Teorema}

\newmdtheoremenv[
  backgroundcolor=gray!10,
  linewidth=0pt,
  innertopmargin=10pt,
  innerbottommargin=10pt,
  innerrightmargin=10pt,
  innerleftmargin=10pt,
  nobreak
]{definition}{Definizione}

\newmdtheoremenv[
  backgroundcolor=gray!10,
  linecolor=myblue,
  linewidth=1.5pt,
  innertopmargin=10pt,
  innerbottommargin=10pt,
  innerrightmargin=10pt,
  innerleftmargin=10pt,
  nobreak
]{proposition}{Proposizione}

% Configurazioni di fancyhdr
\pagestyle{fancy}
\fancyhf{} % Pulisce tutti i campi
\fancyhead[LE,RO]{\thepage} % Numero di pagina
\fancyhead[RE]{\nouppercase{\leftmark}} % Titolo del capitolo
\fancyhead[LO]{\nouppercase{\rightmark}} % Titolo della sezione
\renewcommand{\headrulewidth}{0.5pt}
\renewcommand{\footrulewidth}{0pt}

\geometry{a4paper, left=20mm, right=20mm, top=20mm, bottom=20mm}

% Pacchetto Hyperref (inserito per ultimo)
\usepackage{hyperref}

\begin{document}

\frontmatter
\maketitle
\tableofcontents
\mainmatter


\chapter{Proprietà e calcolo di limiti}

\section{Teorema di unicità del limite e permanenza del segno}

\begin{theorem}[\textbf{Teorema di unicità del limite}]
    Sia \(f(x)\) definita in un intorno di \(x_0\). Se \(\exists \lim_{x \to x_0} f(x) = l\) allora \(l\) è unico.
\end{theorem}


\begin{theorem}[\textbf{Teorema di permanenza del segno}]
    Supponiamo che \(f\) ammetta limite \(L\) (finito o infinito) per \(x\) tendente a \(c\). Se \(L > 0\) oppure \(L = +\infty\), esiste un intorno \(I(c)\) di \(c\) tale che \(f\) è strettamente positiva in \(I(c) \setminus \{c\}\). Un risultato analogo vale per il segno negativo.

    \textbf{Osservazione}
    Vale anche il contrario, se \(f\) è strettamente positiva in un intorno di \(c\), allora \(\lim_{x \to c} f(x) \geq 0\).
\end{theorem}

\begin{theorem}[\textbf{Teorema di Limitatezza Locale}]
    Se una funzione \( f \) ammette limite finito per \( x \to c \), allora esiste un intorno \( I(c) \) di \( c \) e una costante \( M > 0 \) tale che, per ogni \( x \) in \( \text{dom}\,f \cap I(c) \setminus \{c\} \), si ha \( |f(x)| < M \).
\end{theorem}

\subsection*{Considerazione}

I tre teoremi presentati, pur trattando aspetti differenti dei limiti, sono strettamente interconnessi e offrono una visione completa del comportamento di una funzione vicino a un punto di interesse. Il Teorema di unicità del limite assicura che, quando un limite esiste in un punto, questo sia l'unico valore a cui la funzione tende, eliminando ambiguità e fornendo coerenza matematica. Il Teorema di permanenza del segno e il Teorema di Limitatezza Locale, invece, offrono informazioni preziose sul comportamento della funzione nell'intorno del punto: il primo garantisce la coerenza del segno della funzione con il segno del suo limite, mentre il secondo fornisce una garanzia di limitatezza, assicurando che i valori della funzione non diventino arbitrariamente grandi o piccoli. Insieme, questi teoremi formano un quadro potente per comprendere e analizzare il comportamento delle funzioni in prossimità dei loro punti di limite.

\section{Teorema del confronto}

\begin{theorem}[\textbf{Primo Teorema del confronto}]
    Supponiamo che, per \( x \) tendente a \( c \), la funzione \( f \) abbia limite \( l \) mentre la funzione \( g \) abbia limite \( m \) (entrambi finiti o infiniti). Se esiste un intorno \( I(c) \) di \( c \) tale che \( f(x) < g(x) \) in \( I(c) \setminus \{c\} \), allora \( \ell < m \).
\end{theorem}

\begin{theorem}[\textbf{Secondo Teorema del Confronto - Caso Finito}]
    Siano date tre funzioni \(f\), \(g\), \(h\); supponiamo che \(f\) e \(h\) abbiano lo stesso limite finito per \(x\) tendente a \(c\):
    \[\lim_{x \to c} f(x) = \lim_{x \to c} h(x) = L.\]
    Se esiste un intorno \(I(c)\) di \(c\) nel quale siano definite le tre funzioni (tranne al più nel punto \(c\)) e tale che
    \[f(x) \leq g(x) \leq h(x), \forall x \in I(c) \setminus \{c\},\]
    allora si ha anche
    \[\lim_{x \to c} g(x) = L.\]
\end{theorem}
\begin{figure}[H]
    \centering
    \begin{tikzpicture}[
            declare function={
                    h(\x) = -\x^2+4*\x-2);
                    f(\x) = \x^2;
                    g(\x) = 2*\x-1;
                }]
        \begin{axis}
            [
                xlabel={\( x \)},
                ylabel={\( f(x) \)},
                xmin=-2, xmax=3,
                ymin=-2, ymax=3,
                xtick={-2,-1,...,3},
                ytick={-2,-1,...,3},
                grid=both,
                axis lines=middle,
                axis line style=->,
                grid style={line width=.1pt, draw=gray!10},
                major grid style={line width=.2pt,draw=gray!50},
                minor tick num=5,
                enlargelimits={abs=0.5},
                xlabel style={at={(ticklabel* cs:1)}, anchor=north west},
                ylabel style={at={(ticklabel* cs:1)}, anchor=south west},
                samples=100,
                unbounded coords=jump, % Gestisce i punti non definiti
                legend pos=south east
            ]
            \addplot[domain=-3:4,ultra thick, color=red]{h(x)};
            \addplot[domain=-3:3,ultra thick, color=black]{f(x)};
            \addplot[domain=-3:3,ultra thick, color=myblue]{g(x)};
            \filldraw[myblue] (axis cs: 1,1) circle (2pt);
            \draw [dashed, myblue] (axis cs: 1,0) -- (axis cs: 1,1);
            \draw [dashed, myblue] (axis cs: 0,1) -- (axis cs: 1,1);
            \node [font=\small, above right, myblue] at (axis cs: 1,0) {\(x_0\)};
            \node [font=\small, above right, myblue] at (axis cs: 0,1) {\(L\)};
            \legend{\(h(x)\), \(f(x)\), \(g(x)\)}
        \end{axis}
    \end{tikzpicture}
    \caption{Esempio di funzioni che soddisfano le ipotesi del secondo teorema del confronto.}
    \label{fig:secondo-teorema-confronto}
\end{figure}
\subsubsection{Esempio notevole}
Il limite della funzione  \(y = \frac{sinx}{x} \) per \(x \to 0\) non può essere calcolato direttamente.
Tuttavia è possibile dimostrare direttamente il limite
fondamentale
\[\lim_{x \to 0}\frac{\sin x}{x} = 1\]

\begin{enumerate}
    \item Determinazione Diretta del Limite
          Il limite della funzione \( y = \frac{\sin x}{x} \) per \( x \to 0 \) non può essere determinato applicando il Teorema 5.5, in quanto sia il numeratore sia il denominatore tendono a 0. Tuttavia, possiamo dimostrare il limite fondamentale direttamente. Osserviamo che
          \[
              y = \frac{\sin x}{x} \quad \text{e pari, infatti} \quad \frac{\sin(-x)}{-x} = \frac{\sin x}{x}.
          \]
          Dunque, è sufficiente far tendere \( x \) a 0 per valori positivi.

          Ricordando una disuguaglianza trigonometrica fondamentale, per ogni \( x > 0 \) si ha \( \sin x < x \), cioè
          \[
              \frac{\sin x}{x} < 1.
          \]\\

    \item Limitazione Inferiore e Superiore
          Per ottenere una limitazione inferiore, supponiamo \( x < \frac{\pi}{2} \) e consideriamo sulla circonferenza trigonometrica il punto \( A \) di coordinate \( (1,0) \), il punto \( P \) di coordinate \( (\cos x, \sin x) \) e il punto \( Q \) di coordinate \( (1, \tan x) \). Il settore circolare \( OAP \) è strettamente contenuto nel triangolo \( OAQ \), dunque
          \[
              \text{area } OAP < \text{area } OAQ.
          \]
          Essendo
          \[
              \text{area } OAP = \frac{1}{2} x \quad \text{e} \quad \text{area } OAQ = \frac{1}{2} \tan x,
          \]
          si ha
          \[
              \frac{1}{2} x < \frac{1}{2} \tan x, \quad \text{cioè} \quad \cos x < \frac{\sin x}{x}.
          \]
          In conclusione, per \( 0 < x < \frac{\pi}{2} \), abbiamo
          \[
              \cos x < \frac{\sin x}{x} < 1.
          \]
          \begin{figure}[H]
              \centering
              \begin{tikzpicture}[
                      declare function={
                              f(\x) = sqrt(1-\x^2);
                              g(\x) = \x;
                          }]
                  \begin{axis}
                      [
                          xmin=0, xmax=1,
                          ymin=0, ymax=1,
                          xtick=\empty,
                          ytick=\empty,
                          axis equal,
                          grid=both,
                          axis lines=middle,
                          axis line style=->,
                          grid style={line width=.1pt, draw=gray!10},
                          major grid style={line width=.2pt,draw=gray!50},
                          minor tick num=5,
                          enlargelimits={abs=0.5},
                          xlabel style={at={(ticklabel* cs:1)}, anchor=north west},
                          ylabel style={at={(ticklabel* cs:1)}, anchor=south west},
                          samples=100,
                          unbounded coords=jump,
                          legend pos=south east
                      ]
                      \node (A) at (axis cs:0,0) {};
                      \node (B) at (axis cs:1,0) {};
                      \node (C) at (axis cs:1,1) {};
                      \filldraw[fill=lightblue] (A.center) -- (B.center) -- (C.center) -- cycle;
                      \addplot[domain=0:1,ultra thick]{f(x)};
                      \addplot[domain=0:1.2,ultra thick]{g(x)};
                      \draw [ultra thick] (axis cs: 1,0) -- (axis cs: 1,1.2);
                      \node [below right] at (axis cs: 0.3,0.2) {\(x\)};
                      \node [below left] at (axis cs: 0,0) {\(O\)};
                      \node [below] at (axis cs: 1,0) {\(A\)};
                      \node [above] at (axis cs: 0.7,0.71) {\(P\)};
                      \node [above left] at (axis cs: 1,1) {\(Q\)};
                      \draw [dashed, myblue] (axis cs: 0.7,0.71) -- (axis cs: 0.7,0);
                      \draw [dashed, myblue] (axis cs: 0.7,0.71) -- (axis cs: 0,0.71);
                      \node [above, myblue] at (axis cs: 0.35,0.71) {\(\cos x\)};
                      \node [right, myblue] at (axis cs: 0.7,0.35) {\(\sin x\)};
                  \end{axis}
              \end{tikzpicture}
          \end{figure}
          La continuità della funzione coseno assicura che \( \lim_{x \to 0^+} \cos x = 1 \). Pertanto, applicando il secondo Teorema del Confronto, otteniamo la tesi.
    \item Studio del Comportamento Limite per \( x \to +\infty \)
          Vogliamo ora studiare il comportamento limite della funzione \( g(x) = \frac{\sin x}{x} \) per \( x \) tendente a \( +\infty \). Ricordiamo che per ogni \( x \) reale si ha
          \[
              -1 \leq \sin x \leq 1.
          \]
          Dividendo ciascun termine per \( x > 0 \), le disuguaglianze si conservano; pertanto, in ogni intorno \( I(c) \) di \( +\infty \) si ha
          \[
              -\frac{1}{x} \leq \frac{\sin x}{x} \leq \frac{1}{x}.
          \]
          Posto \( f(x) = -\frac{1}{x} \), \( h(x) = \frac{1}{x} \) e osservato che \( \lim_{x \to +\infty} \frac{1}{x} = 0 \), deduciamo dal teorema precedente che
          \[
              0 \leq \frac{\sin x}{x} \leq 0.
          \]
          \[
              \lim_{x \to +\infty} \frac{\sin x}{x} = 0.
          \]
          \begin{figure}[H]
              \centering
              \begin{tikzpicture}[
                      declare function={
                              h(\x) = 1/\x;
                              f(\x) = sin(deg(\x))/\x;
                              g(\x) = -1/\x;
                          }]
                  \begin{axis}
                      [
                          xlabel={\( x \)},
                          ylabel={\( f(x) \)},
                          xmin=-5, xmax=5,
                          ymin=-5, ymax=5,
                          grid=both,
                          axis lines=middle,
                          axis line style=->,
                          grid style={line width=.1pt, draw=gray!10},
                          major grid style={line width=.2pt,draw=gray!50},
                          minor tick num=3,
                          enlargelimits={abs=0.5},
                          xlabel style={at={(ticklabel* cs:1)}, anchor=north west},
                          ylabel style={at={(ticklabel* cs:1)}, anchor=south west},
                          samples=100,
                          unbounded coords=jump, % Gestisce i punti non definiti
                          legend pos=south east
                      ]
                      \addplot[domain=-5:5,ultra thick, color=red] {ifthenelse(abs(x)>0.1, h(x), NaN)};
                      \addplot[domain=-5:5,ultra thick, color=blue] {ifthenelse(abs(x)>0.1, g(x), NaN)};
                      \addplot[domain=-5:5,ultra thick, color=myblue]{f(x)};
                      \legend{\(h(x)\), \(h(x)\), \(f(x)\), \(g(x)\), \(g(x)\)}
                  \end{axis}
              \end{tikzpicture}
          \end{figure}
\end{enumerate}
\begin{theorem}[\textbf{Secondo Teorema del Confronto - Caso Infinito}]
    Siano date due funzioni \(f\) e \(g\) ed esista il limite
    \[\lim_{x \to c} f(x) = +\infty.\]
    Se esiste un intorno \(I(c)\) di \(c\) nel quale siano definite entrambe le funzioni (tranne al più nel punto \(c\)) e tale che
    \[f(x) \leq g(x), \forall x \in I(c) \setminus \{c\},\]
    allora si ha anche
    \[\lim_{x \to c} g(x) = +\infty.\]
    Un risultato analogo vale nel caso del limite \(-\infty\).
\end{theorem}
\subsection*{Considerazione sui Teoremi del Confronto}

I Teoremi del Confronto, nella loro varietà, offrono una visione profonda e versatile nel campo dell'analisi dei limiti. Questi teoremi sono essenziali per determinare il comportamento asintotico delle funzioni in punti critici, senza la necessità di calcolare direttamente i limiti, che possono talvolta essere complessi o irrisolvibili con metodi standard.

Il \textbf{Primo Teorema del Confronto} stabilisce una relazione diretta tra i limiti di due funzioni: se una funzione è sistematicamente inferiore a un'altra in un intorno di un punto e entrambe tendono a un limite, allora il limite della funzione inferiore è al massimo uguale a quello della funzione superiore. Questo principio è fondamentale per comprendere come il comportamento di una funzione sia influenzato dalla sua "posizione" relativa rispetto ad altre funzioni.

Il \textbf{Secondo Teorema del Confronto - Caso Finito} amplia questa idea al caso di tre funzioni, permettendo di dedurre il limite di una funzione intermedia quando è "schiacciata" tra due funzioni con lo stesso limite noto. Questo teorema è particolarmente utile per stabilire limiti in situazioni in cui il calcolo diretto è complesso, utilizzando invece funzioni ausiliarie più gestibili.

Infine, il \textbf{Secondo Teorema del Confronto - Caso Infinito} si occupa di situazioni in cui almeno una delle funzioni tende all'infinito. Questo teorema è cruciale per stabilire il comportamento asintotico delle funzioni che esibiscono crescita o decrescita illimitate, fornendo un metodo per trattare limiti altrimenti difficili da affrontare.

In conclusione, i Teoremi del Confronto sono strumenti indispensabili nell'analisi dei limiti, fornendo un metodo potente e flessibile per esplorare e comprendere il comportamento delle funzioni in prossimità di punti critici. La loro applicazione spazia dall'analisi reale all'ingegneria, dimostrando ancora una volta l'importanza fondamentale dei concetti di limite e confronto nell'ambito matematico.

\section{Algebra dei Limiti}
\begin{proposition}[\textbf{Operazioni definite}]
    \begin{align*}
        +\infty + s                      & = +\infty \quad   & \text{(se } s \in \mathbb{R} \text{ oppure } s = +\infty\text{)}                   \\
        -\infty + s                      & = -\infty \quad   & \text{(se } s \in \mathbb{R} \text{ oppure } s = -\infty\text{)}                   \\
        \pm\infty \cdot s                & = \pm\infty \quad & \text{(se } s > 0 \text{ oppure } s = +\infty\text{)}                              \\
        \pm\infty \cdot s                & = \mp\infty \quad & \text{(se } s < 0 \text{ oppure } s = -\infty\text{)}                              \\
        \displaystyle\frac{\pm\infty}{s} & = \pm\infty \quad & \text{(se } s > 0\text{)}                                                          \\
        \displaystyle\frac{\pm\infty}{s} & = \mp\infty \quad & \text{(se } s < 0\text{)}                                                          \\
        \displaystyle\frac{s}{0}         & = \infty \quad    & \text{(se } s \in \mathbb{R} \setminus \{0\} \text{ oppure } s = \pm\infty\text{)} \\
        \displaystyle\frac{s}{\pm\infty} & = 0 \quad         & \text{(se } s \in \mathbb{R}\text{)}
    \end{align*}
\end{proposition}
\begin{proposition}[\textbf{Operazioni indefinite}]
    \begin{align*}
        \infty - \infty       & = \text{indeterminato} \\
        \frac{\infty}{\infty} & = \text{indeterminato} \\
        \frac{0}{0}           & = \text{indeterminato} \\
        0 \cdot \infty        & = \text{indeterminato}
    \end{align*}
\end{proposition}
\begin{theorem}[\textbf{Algebra dei limiti}]
    Supponiamo che, per \( x \) tendente a \( c \), la funzione \( f \) ammetta limite \( \ell \) (finito o infinito) e la funzione \( g \) ammetta limite \( m \) (anch'esso finito o infinito). Allora, ogniqualvolta l'espressione al secondo membro è definita, si ha
    \begin{align*}
        \lim_{x \to c} (f(x) \pm g(x))   & = \ell \pm m,    \\
        \lim_{x \to c} (f(x)g(x))        & = \ell m,        \\
        \lim_{x \to c} \frac{f(x)}{g(x)} & = \frac{\ell}{m}
    \end{align*}
    \footnote[1]{in quest'ultimo caso supponiamo inoltre che \(g(x) \neq 0\) in un intorno di \(c\) escluso al più il punto \(c\).}
\end{theorem}
\noindent\begin{minipage}{\textwidth}
    \begin{proposition}[\textbf{Limiti di funzioni elementari}]
        i limiti più significativi delle principali funzioni elementari:
        \begin{center}
            \renewcommand{\arraystretch}{1.5} % Aumenta lo spazio tra le righe
            \begin{tabular}{l l}
                $\displaystyle \lim_{x \to +\infty} x^a = 0$                                                                                                & se $a > 0$                                                                    \\
                $\displaystyle \lim_{x \to 0^+} x^a = 0$, $\displaystyle \lim_{x \to 0^+} x^a = +\infty$                                                    & se $a < 0$                                                                    \\
                \\
                $\displaystyle \lim_{x \to +\infty} \frac{a_n x^n + \ldots + a_1 x + a_0}{b_m x^m + \ldots + b_1 x + b_0}$                                  &                                                                               \\
                $= \infty$                                                                                                                                  & se $n > m$                                                                    \\
                $= \frac{a_n}{b_m}$                                                                                                                         & se $n = m$                                                                    \\
                $= 0$                                                                                                                                       & se $n < m$                                                                    \\
                \\
                $\displaystyle \lim_{x \to +\infty} a^x = +\infty$, $\displaystyle \lim_{x \to +\infty} a^x = 0$                                            & se $a > 1$                                                                    \\
                $\displaystyle \lim_{x \to +\infty} a^x = 0$, $\displaystyle \lim_{x \to +\infty} a^x = +\infty$                                            & se $a < 1$                                                                    \\
                \\
                $\displaystyle \lim_{x \to 0^+} \log_a x = +\infty$, $\displaystyle \lim_{x \to 0^+} \log_a x = -\infty$                                    & se $a > 1$                                                                    \\
                $\displaystyle \lim_{x \to 0^+} \log_a x = -\infty$, $\displaystyle \lim_{x \to 0^+} \log_a x = +\infty$                                    & se $a < 1$                                                                    \\
                \\
                $\displaystyle \lim_{x \to \pm\infty} \sin x$, $\displaystyle \lim_{x \to \pm\infty} \cos x$, $\displaystyle \lim_{x \to \pm\infty} \tan x$ & non esistono                                                                  \\
                $\displaystyle \lim_{x \to (k + \frac{1}{2})\pi \pm 0} \tan x$                                                                              & $= \pm\infty, \forall k \in \mathbb{Z}$                                       \\
                \\
                $\displaystyle \lim_{x \to \pm 1} \arcsin x$                                                                                                & $= \pm \frac{\pi}{2} = \arcsin(\pm 1)$                                        \\
                $\displaystyle \lim_{x \to +1} \arccos x = 0$                                                                                               & $= \arccos 1$, $\displaystyle \lim_{x \to -1} \arccos x = \pi = \arccos (-1)$ \\
                $\displaystyle \lim_{x \to \pm\infty} \arctan x$                                                                                            & $= \pm \frac{\pi}{2}$                                                         \\
            \end{tabular}
        \end{center}
    \end{proposition}
\end{minipage}
\section{Teorema di sostiuzione}
\noindent\begin{minipage}{\textwidth}
    \begin{theorem}[\textbf{Teorema di Sostituzione}]
        Supponiamo che esista (finito o infinito)
        \[\lim_{x \to c} f(x) = \ell.\]
        Sia poi \( g \) una funzione definita in un intorno di \( \ell \), (escluso al più il punto \( \ell \)) e tale che:
        \begin{enumerate}
            \item se \( \ell \in \mathbb{R} \), \( g \) è continua in \( \ell \);
            \item se \( \ell = +\infty \) oppure \( \ell = -\infty \), esiste (finito o infinito) \(\lim_{y \to \ell} g(y)\).
        \end{enumerate}
        Allora, esiste il limite per \( x \) tendente a \( c \) della funzione composta \( g \circ f \) e si ha
        \[\lim_{x \to c} (g \circ f)(x) = \lim_{y \to \ell} g(y).\]
    \end{theorem}
\end{minipage}
\subsection*{Considerazione}
In pratica, nel calcolo del limite della funzione \( g(f(x)) \), il Teorema di Sostituzione ci permette di effettuare la sostituzione \( y = f(x) \) e successivamente studiare il comportamento limite della funzione \( g(y) \). Questo approccio semplifica l'analisi, in particolare quando la funzione composta \( g(f(x)) \) presenta una struttura complicata o non immediatamente evidente. Analizzando il comportamento di \( g(y) \) mentre \( y \) tende al limite di \( f(x) \), possiamo determinare il limite della funzione composta.

\section{Limiti notevoli}

\begin{proposition}
    Vale il seguente risultato
    \[
        \lim_{{x \to \infty}} \left(1 + \frac{1}{x}\right)^x = e
    \]
\end{proposition}

\noindent\begin{minipage}{\textwidth}
    \begin{proposition}[\textbf{Limiti Notevoli}]
        I limiti notevoli delle principali funzioni elementari:
        \begin{center}
            \renewcommand{\arraystretch}{2} % Aumenta lo spazio tra le righe
            \begin{tabular}{l l}
                $\displaystyle \lim_{x \to 0} \frac{\sin(x)}{x} = 1$                       &                               \\
                $\displaystyle \lim_{x \to 0} \frac{\cos(x)}{x} = \text{indeterminato}$    &                               \\
                $\displaystyle \lim_{x \to \infty} \frac{\sin(x)}{x} = 0$                  &                               \\
                $\displaystyle \lim_{x \to \infty} \frac{\cos(x)}{x} = 0$                  &                               \\
                $\displaystyle \lim_{x \to 0} \frac{\tan(x)}{x} = 1$                       &                               \\
                $\displaystyle \lim_{x \to 0} \frac{1 - \cos(x)}{x^2} = \frac{1}{2}$       &                               \\
                $\displaystyle \lim_{x \to \infty} \left(1 + \frac{a}{x}\right)^x = e^a$   & con $a \in \mathbb{R}$        \\
                $\displaystyle \lim_{x \to 0} \left(1 + x\right)^{\frac{1}{x}} = e$        &                               \\
                $\displaystyle \lim_{x \to 0} \frac{\log_a(1 + x)}{x} = \frac{1}{\log(a)}$ & con $a > 0$ e $a \neq 1$      \\
                $\displaystyle \lim_{x \to 0} \frac{\log(1 + x)}{x} = 1$                   & (logaritmo naturale)          \\
                $\displaystyle \lim_{x \to 0} \frac{a^x - 1}{x} = \log(a)$                 & con $a > 0$                   \\
                $\displaystyle \lim_{x \to 0} \frac{e^x - 1}{x} = 1$                       & (base dei logaritmi naturali) \\
                $\displaystyle \lim_{x \to 0} \frac{(1 + x)^\alpha - 1}{x} = \alpha$       & con $\alpha \in \mathbb{R}$   \\
            \end{tabular}
        \end{center}
    \end{proposition}
\end{minipage}

\section{Forme indeterminate di tipo esponenziale}

\begin{proposition}[Identità Fondamentale]
    Per ogni $x > 0$, abbiamo che
    \[
        f(x) = x = e^{\log(x)}
    \]
    Questo perché $\log(x)$ è il logaritmo naturale di $x$, e $e$ è la base del logaritmo naturale, quindi le due operazioni si annullano a vicenda.
\end{proposition}
\subsection{Analisi dei Limiti con Funzioni Esponenziali e Logaritmiche}
\noindent\begin{minipage}{\textwidth}
    \begin{proposition}[\textbf{Trasformazione di una funzione in una forma esponenziale}]
        Supponiamo di avere una funzione esponenziale che vogliamo analizzare.
        \[
            \displaystyle h(x) = [f(x)]^{g(x)}
        \]
        Possiamo trasformarla in una forma che coinvolge esponenti e logaritmi, ovvero
        \[
            \displaystyle h(x) = e^{g(x) \log(f(x))}
        \]
        dove $g(x)$ e $f(x)$ sono altre funzioni che dipendono da $x$.
    \end{proposition}
\end{minipage}

\begin{proposition}[Trasformazione di una funzione in forma esponenziale per calcolo limit]
    Vale la seguente proprietà
    \begin{align*}
        \renewcommand{\arraystretch}{2}
        \displaystyle \lim_{x \to c} h(x) & = \lim_{x \to c} [f(x)]^{g(x)}                      \\
        \displaystyle                     & = \lim_{x \to c} e^{g(x) \log(f(x))}                \\
        \displaystyle                     & = \exp\left[{\lim_{x \to c} g(x) \log(f(x))}\right] \\
        \displaystyle                     & = e^{\lim_{x \to c} g(x) \log(f(x))}                \\
    \end{align*}
\end{proposition}
\begin{proposition}[Forme indeterminate connesse a funzioni Esponenziali]
    \[
        0 \cdot \infty
    \]
    \[
        0^0
    \]
    \[
        \infty^0
    \]
\end{proposition}

\end{document}